%% paper47_darstellungstheorie_de.tex
%% Paper 47 (DE): Darstellungstheorie — Von Gruppenalgebren zum Langlands-Programm
%% Autor: Michael Fuhrmann
%% Projekt: Specialist Mathematics, Build 135
%% Datum: 2026-03-12

\documentclass[12pt,a4paper]{amsart}
\usepackage{amsmath,amssymb,amsthm}
\usepackage{mathtools}
\usepackage[utf8]{inputenc}
\usepackage[T1]{fontenc}
\usepackage{hyperref}
\usepackage{enumitem,booktabs,array}
\usepackage{geometry}
usepackage{ytableau}
\geometry{margin=2.5cm}

\newtheorem{theorem}{Satz}[section]
\newtheorem{lemma}[theorem]{Lemma}
\newtheorem{korollar}[theorem]{Korollar}
\newtheorem{proposition}[theorem]{Proposition}
\newtheorem{vermutung}[theorem]{Vermutung}
\theoremstyle{definition}
\newtheorem{definition}[theorem]{Definition}
\newtheorem{beispiel}[theorem]{Beispiel}
\newtheorem{bemerkung}[theorem]{Bemerkung}

\title{Darstellungstheorie:\\
Von Gruppenalgebren zum Langlands-Programm}

\author{Michael Fuhrmann}
\address{Specialist Mathematics Project, Build 135}
\email{specialist-maths@localhost}
\date{\today}

\begin{document}

\begin{abstract}
Die Darstellungstheorie stellt einen systematischen Rahmen bereit, um abstrakte
Gruppen durch ihre konkreten Wirkungen auf Vektorräumen zu verstehen.  Ausgehend
von der grundlegenden Definition eines Gruppenhomomorphismus in allgemeine lineare
Gruppen entwickeln wir den Kern des Faches: Schurs Lemma, den Maschke-Satz über
vollständige Reduzibilität sowie die reiche Orthogonalitätstheorie der Charaktere.
Anschließend untersuchen wir die Darstellungstheorie der symmetrischen Gruppen
$S_n$ mittels Young-Tableaux und Specht-Moduln und wechseln zum stetigen Rahmen
mit der Höchstgewicht-Klassifikation der Darstellungen von $\mathrm{GL}_n$ und
$\mathrm{SL}_n$.  Anwendungen reichen von Burnsides $pq$-Satz in der reinen
Gruppentheorie bis zur molekularen Symmetrie und Kristallographie.  Abschließend
geben wir eine Einführung in Galois-Darstellungen und das Langlands-Programm,
eines der tiefsten Vermutungsgebäude der modernen Mathematik.
\end{abstract}

\maketitle

\tableofcontents

%% ============================================================
\section{Einleitung}
%% ============================================================

Darstellungstheorie ist die Kunst, abstrakte algebraische Strukturen dadurch zu
verstehen, dass man sie auf Vektorräumen wirken lässt.  Eine \emph{Darstellung}
einer Gruppe $G$ auf einem Vektorraum $V$ ist schlicht ein Gruppenhomomorphismus
\[
    \rho \colon G \longrightarrow \mathrm{GL}(V),
\]
wobei $\mathrm{GL}(V)$ die Gruppe der invertierbaren linearen Endomorphismen von
$V$ bezeichnet.  Durch Wahl einer Basis identifiziert man $\mathrm{GL}(V) \cong
\mathrm{GL}(n,k)$ für $n = \dim_k V$, so dass $\rho$ jedem Gruppenelement $g \in G$
eine invertierbare $n \times n$-Matrix $\rho(g)$ zuordnet.

Das Fach steht an einer bemerkenswerten Kreuzung von Algebra und linearer
Algebra.  Auf der einen Seite übersetzt es die oft schwer handhabbare Kombinatorik
einer Gruppe in die gut verstandene Geometrie der Matrixmultiplikation.  Auf der
anderen Seite liefert es das algebraische Fundament großer Teile der modernen
Physik: Die Darstellungen der Rotationsgruppe $\mathrm{SO}(3)$ und ihrer
zweifachen Überlagerung $\mathrm{SU}(2)$ beschreiben den quantenmechanischen
Spin, während die Darstellungen der Lorentz-Gruppe und von $\mathrm{SU}(3)$
Elementarteilchen im Standardmodell klassifizieren.

Aus rein mathematischer Perspektive hat das Fach enorme Fortschritte angestoßen.
Burnsides Beweis von 1904, dass jede Gruppe der Ordnung $p^a q^b$ ($p$, $q$
prim) auflösbar ist, stützt sich in wesentlicher Weise auf Charaktertheorie,
und bis heute gibt es keinen rein gruppentheoretischen Beweis.  Die Klassifikation
der endlichen einfachen Gruppen wäre ohne die eingehende Untersuchung modularer
Darstellungen nicht denkbar.  Und an der Forschungsfront sagt das
Langlands-Programm---von Robert Langlands 1967 in einem berühmten Brief an
André Weil formuliert---tiefe Korrespondenzen zwischen automorphen Darstellungen
reduktiver Gruppen und Galois-Darstellungen voraus und vereinigt weite Teile der
Zahlentheorie, Geometrie und Analysis.

\medskip

\noindent\textbf{Überblick.} Abschnitt~\ref{sec:grund} führt die
Grunddefinitionen und die wichtigsten Beispielklassen ein.
Abschnitt~\ref{sec:schur} beweist Schurs Lemma und zieht seine fundamentalen
Konsequenzen.  Abschnitt~\ref{sec:maschke} beweist die vollständige
Reduzibilität über Körpern, deren Charakteristik die Gruppenordnung nicht
teilt.  Abschnitt~\ref{sec:charaktere} entwickelt die Charaktertheorie
inklusive der beiden Orthogonalitätsrelationen.  Abschnitt~\ref{sec:sn}
behandelt die Darstellungstheorie der symmetrischen Gruppe mittels
Young-Tableaux.  Abschnitt~\ref{sec:gln} skizziert die Höchstgewicht-Theorie
für $\mathrm{GL}_n$ und $\mathrm{SL}_n$.  Abschnitt~\ref{sec:anwendungen}
sammelt drei klassische Anwendungen.  Abschnitt~\ref{sec:langlands} führt in
Galois-Darstellungen und das Langlands-Programm ein.

%% ============================================================
\section{Grundbegriffe der Darstellungstheorie}
\label{sec:grund}
%% ============================================================

Sei $G$ eine Gruppe und $k$ ein Körper.  Sofern nicht anders angegeben,
sei $k = \mathbb{C}$.

\begin{definition}
Eine \emph{Darstellung} von $G$ über $k$ ist ein Paar $(V,\rho)$, wobei $V$ ein
$k$-Vektorraum und $\rho \colon G \to \mathrm{GL}(V)$ ein Gruppenhomomorphismus
ist.  Die \emph{Dimension} (oder der \emph{Grad}) von $(V,\rho)$ ist
$\dim_k V$.
\end{definition}

Äquivalent ist eine Darstellung ein $k[G]$-Modul, wobei $k[G]$ die
\emph{Gruppenalgebra} von $G$ über $k$ bezeichnet: den freien $k$-Vektorraum auf
den Elementen von $G$ mit Multiplikation, die diejenige von $G$ linear fortsetzt.

\begin{beispiel}[Triviale Darstellung]
Die \emph{triviale Darstellung} ist $\rho_{\mathrm{triv}} \colon G \to
\mathrm{GL}(k)$, $g \mapsto 1$.  Jedes Element wirkt als Identität.
\end{beispiel}

\begin{beispiel}[Reguläre Darstellung]
Die \emph{reguläre Darstellung} $(\mathbb{C}[G], \rho_{\mathrm{reg}})$ ist durch
Links-Multiplikation definiert:
\[
    \rho_{\mathrm{reg}}(g)(e_h) = e_{gh},
\]
wobei $\{e_h\}_{h \in G}$ die Standardbasis von $\mathbb{C}[G]$ ist.  Ihre
Dimension ist $|G|$ und sie enthält jede irreduzible Darstellung von $G$ mit einer
Multiplizität gleich ihrer eigenen Dimension.
\end{beispiel}

\begin{beispiel}[Permutationsdarstellung]
Wirkt $G$ auf einer endlichen Menge $X = \{x_1,\ldots,x_n\}$, so ist die
\emph{Permutationsdarstellung} $V = \mathbb{C}^X$ durch
$\rho(g)(e_{x_i}) = e_{g \cdot x_i}$ definiert.  Die reguläre Darstellung ist
der Spezialfall $X = G$ mit der Links-Multiplikationswirkung.
\end{beispiel}

\begin{beispiel}[Vorzeichen-Darstellung von $S_n$]
Die symmetrische Gruppe $S_n$ besitzt eine eindimensionale Darstellung
$\mathrm{sgn} \colon S_n \to \{+1,-1\}$, die jede Permutation auf ihr Vorzeichen
abbildet.  Für $n \geq 2$ ist sie von der trivialen Darstellung verschieden und
liefert unser erstes Beispiel einer nicht-trivialen irreduziblen Darstellung.
\end{beispiel}

\begin{definition}
Eine \emph{Unterdarstellung} von $(V,\rho)$ ist ein $G$-stabiler Unterraum
$W \subseteq V$, d.h.\ $\rho(g)W \subseteq W$ für alle $g \in G$.  Die
Darstellung $(V,\rho)$ heißt \emph{irreduzibel}, wenn ihre einzigen
Unterdarstellungen $\{0\}$ und $V$ selbst sind.
\end{definition}

\begin{definition}
Ein \emph{Morphismus} (oder \emph{$G$-äquivariante Abbildung}, oder
\emph{Intertwiner}) zwischen Darstellungen $(V,\rho)$ und $(W,\sigma)$ ist
eine lineare Abbildung $T \colon V \to W$ mit
\[
    T \circ \rho(g) = \sigma(g) \circ T \quad \text{für alle } g \in G.
\]
Zwei Darstellungen heißen \emph{isomorph} (oder \emph{äquivalent}), wenn
ein bijektiver Intertwiner zwischen ihnen existiert.
\end{definition}

\begin{beispiel}[Darstellungen von $\mathbb{Z}/2\mathbb{Z}$]
Die Gruppe $\mathbb{Z}/2\mathbb{Z} = \{0,1\}$ besitzt genau zwei irreduzible
komplexe Darstellungen, beide eindimensional:
\begin{itemize}
  \item Triviale Darstellung: $\rho(0) = (1)$, $\rho(1) = (1)$.
  \item Vorzeichen-Darstellung: $\rho(0) = (1)$, $\rho(1) = (-1)$.
\end{itemize}
Dies sind alle irreduziblen Darstellungen, da deren Anzahl gleich der Anzahl der
Konjugationsklassen ist, welche für eine abelsche Gruppe der Ordnung $2$ gleich
$2$ ist.
\end{beispiel}

\begin{beispiel}[Darstellungen von $S_3$]
Die Gruppe $S_3$ hat Ordnung $6$ und drei Konjugationsklassen: $\{e\}$,
$\{(12),(13),(23)\}$, $\{(123),(132)\}$.  Damit gibt es genau drei irreduzible
Darstellungen:
\begin{itemize}
  \item Triviale Darstellung (Dimension $1$): $\rho(\sigma) = 1$ für alle $\sigma$.
  \item Vorzeichen-Darstellung (Dimension $1$): $\rho(\sigma) = \mathrm{sgn}(\sigma)$.
  \item Standard-Darstellung (Dimension $2$): orthogonale Wirkung auf das
        2-dimensionale Komplement des Fixpunktraums in der Permutationsdarstellung.
\end{itemize}
Man verifiziert: $1^2 + 1^2 + 2^2 = 6 = |S_3|$.
\end{beispiel}

%% ============================================================
\section{Schurs Lemma}
\label{sec:schur}
%% ============================================================

Schurs Lemma ist der Eckpfeiler aller weiteren Strukturtheorie.  Es besagt, dass
irreduzible Darstellungen so wenige nicht-triviale Intertwiner wie möglich besitzen.

\begin{theorem}[Schurs Lemma]
\label{satz:schur}
Seien $(V,\rho)$ und $(W,\sigma)$ irreduzible Darstellungen von $G$ über einem
Körper $k$, und sei $T \colon V \to W$ eine $G$-äquivariante lineare Abbildung.
\begin{enumerate}[label=\textup{(\roman*)}]
  \item Entweder ist $T = 0$ oder $T$ ist ein Isomorphismus.
  \item Ist $k$ algebraisch abgeschlossen und $V = W$ (mit $\rho = \sigma$), so
        gilt $T = \lambda \cdot \mathrm{id}_V$ für ein Skalar $\lambda \in k$.
\end{enumerate}
\end{theorem}

\begin{proof}
(i) Der Kern $\ker T$ ist eine Unterdarstellung von $V$ und das Bild
$\mathrm{im}\, T$ ist eine Unterdarstellung von $W$.  Da $V$ irreduzibel ist,
gilt $\ker T \in \{\{0\},V\}$.  Da $W$ irreduzibel ist, gilt
$\mathrm{im}\, T \in \{\{0\},W\}$.  Ist $\ker T = V$, so ist $T = 0$.
Andernfalls ist $\ker T = \{0\}$ und $\mathrm{im}\, T = W$, also ist $T$
bijektiv.

(ii) Da $k$ algebraisch abgeschlossen ist, besitzt $T$ einen Eigenwert
$\lambda \in k$.  Die Abbildung $T - \lambda\,\mathrm{id}_V$ ist wieder
$G$-äquivariant und hat einen nicht-trivialen Kern, also ist sie nach Teil (i)
die Nullabbildung, d.h.\ $T = \lambda\,\mathrm{id}_V$.
\end{proof}

\begin{korollar}
\label{kor:schur}
Sei $G$ eine endliche Gruppe und $k = \mathbb{C}$.  Sind $(V,\rho)$ und
$(W,\sigma)$ nicht-isomorphe irreduzible Darstellungen, so ist jede
$G$-äquivariante Abbildung $T \colon V \to W$ die Nullabbildung.  Insbesondere:
\[
    \mathrm{Hom}_G(V,W) = \begin{cases} \mathbb{C} & \text{falls } V \cong W, \\ 0 & \text{sonst.} \end{cases}
\]
\end{korollar}

\begin{korollar}
Jede irreduzible komplexe Darstellung einer abelschen Gruppe ist eindimensional.
\end{korollar}

\begin{proof}
Für jedes $h \in G$ ist die Abbildung $\rho(h) \colon V \to V$ $G$-äquivariant,
da $G$ abelsch ist.  Nach Schurs Lemma gilt $\rho(h) = \lambda_h \cdot
\mathrm{id}_V$.  Damit ist jeder Unterraum von $V$ $G$-stabil.  Die
Irreduzibilität erzwingt $\dim V = 1$.
\end{proof}

\begin{bemerkung}
Die Voraussetzung der algebraischen Abgeschlossenheit in Teil (ii) ist notwendig.
Über $\mathbb{R}$ ist die zweidimensionale Darstellung von $\mathrm{SO}(2)$ durch
Rotationsmatrizen irreduzibel, doch ihr Endomorphismenring ist isomorph zu
$\mathbb{C}$, nicht zu $\mathbb{R}$.
\end{bemerkung}

Die Bedeutung von Schurs Lemma lässt sich an einem weiteren Korollar ablesen,
das der Charaktertheorie zugrunde liegt:

\begin{korollar}[Schur-Orthogonalitätsrelationen, elementare Form]
Seien $\rho \colon G \to \mathrm{GL}(V)$ und $\sigma \colon G \to \mathrm{GL}(W)$
irreduzible Darstellungen.  Dann gilt für alle $i, j, k, l$:
\[
    \frac{1}{|G|} \sum_{g \in G} \rho(g)_{ij}\,\overline{\sigma(g)_{kl}}
    = \frac{\delta_{\rho\sigma}\,\delta_{ik}\,\delta_{jl}}{\dim V}.
\]
\end{korollar}

%% ============================================================
\section{Vollständige Reduzibilität: Der Maschke-Satz}
\label{sec:maschke}
%% ============================================================

\begin{theorem}[Maschke-Satz]
\label{satz:maschke}
Sei $G$ eine endliche Gruppe und $k$ ein Körper mit $\mathrm{char}(k) = 0$ oder
$\mathrm{char}(k) \nmid |G|$.  Dann ist jede Darstellung von $G$ über $k$
\emph{vollständig reduzibel}: Sie zerfällt in eine direkte Summe irreduzibler
Unterdarstellungen.
\end{theorem}

\begin{proof}
Es genügt zu zeigen, dass jede Unterdarstellung ein $G$-stabiles Komplement
besitzt.  Sei $(V,\rho)$ eine Darstellung und $W \subseteq V$ eine
Unterdarstellung.  Da $V$ ein endlich-dimensionaler $k$-Vektorraum ist, existiert
eine lineare Projektion $\pi_0 \colon V \to W$ (mit $\pi_0|_W = \mathrm{id}_W$),
die aber im Allgemeinen nicht $G$-äquivariant ist.  Wir \emph{mitteln} über die
Gruppe, um eine $G$-äquivariante Projektion zu erzeugen:
\[
    \pi \coloneqq \frac{1}{|G|} \sum_{g \in G} \rho(g) \circ \pi_0 \circ \rho(g)^{-1}.
\]
Der Faktor $1/|G|$ ist wohldefiniert, da $|G|$ nach Voraussetzung in $k$
invertierbar ist.  Man verifiziert:
\begin{itemize}
  \item $\pi$ ist $G$-äquivariant: $\rho(h) \circ \pi = \pi \circ \rho(h)$ für alle $h \in G$.
  \item $\pi|_W = \mathrm{id}_W$: Für $w \in W$ und beliebiges $g$ gilt
        $\rho(g)^{-1}w \in W$ (da $W$ $G$-stabil ist), also
        $\pi_0(\rho(g)^{-1}w) = \rho(g)^{-1}w$ und $\rho(g)\rho(g)^{-1}w = w$.
  \item Insbesondere $\pi \circ \pi = \pi$ (Idempotenz).
\end{itemize}
Daher ist $V = W \oplus \ker\pi$ eine Zerlegung in $G$-stabile Unterräume.  Durch
Noethersche Induktion über $\dim V$ terminiert der Prozess und liefert eine
Zerlegung in Irreduzible.
\end{proof}

\begin{bemerkung}
Die Voraussetzung $\mathrm{char}(k) \nmid |G|$ ist scharf.  Über einem Körper der
Charakteristik $p$ mit $p \mid |G|$ können unzerlegbare, aber reduzible
Darstellungen auftreten.  Die \emph{modulare Darstellungstheorie} endlicher
Gruppen (Brauer-Theorie) ist das Studium dieser subtileren Situation.
\end{bemerkung}

Unter den Voraussetzungen des Maschke-Satzes zerfällt jede Darstellung $(V,\rho)$
als
\[
    V \cong \bigoplus_{i} V_i^{\oplus m_i},
\]
wobei die $V_i$ paarweise nicht-isomorphe irreduzible Darstellungen sind und
$m_i \geq 0$ eindeutig bestimmte nicht-negative ganze Zahlen sind, die
\emph{Multiplizitäten} genannt werden.  Die Multiplizitäten lassen sich mithilfe
der im nächsten Abschnitt entwickelten Charaktertheorie berechnen.

%% ============================================================
\section{Charaktere}
\label{sec:charaktere}
%% ============================================================

\begin{definition}
Der \emph{Charakter} einer Darstellung $(V,\rho)$ ist die Funktion
\[
    \chi_V \colon G \longrightarrow k, \quad \chi_V(g) = \mathrm{Spur}(\rho(g)).
\]
\end{definition}

Charaktere sind \emph{Klassenfunktionen}: sie sind konstant auf Konjugationsklassen,
da $\mathrm{Spur}(ABA^{-1}) = \mathrm{Spur}(B)$.  Der Vektorraum der
Klassenfunktionen auf $G$ wird mit $\mathrm{Kl}(G)$ bezeichnet; seine Dimension
ist gleich der Anzahl $r$ der Konjugationsklassen.

\begin{proposition}
\label{prop:char}
Seien $(V,\rho)$ und $(W,\sigma)$ Darstellungen.  Dann:
\begin{enumerate}[label=\textup{(\roman*)}]
  \item $\chi_V(e) = \dim V$.
  \item $\chi_V(g^{-1}) = \overline{\chi_V(g)}$ für $|G| < \infty$, $k = \mathbb{C}$.
  \item $\chi_{V \oplus W} = \chi_V + \chi_W$.
  \item $\chi_{V \otimes W} = \chi_V \cdot \chi_W$.
\end{enumerate}
\end{proposition}

\begin{definition}[Inneres Produkt auf Klassenfunktionen]
Für $k = \mathbb{C}$ ist das \emph{innere Produkt} auf Klassenfunktionen:
\[
    \langle \alpha, \beta \rangle_G \coloneqq \frac{1}{|G|} \sum_{g \in G} \alpha(g)\,\overline{\beta(g)}.
\]
\end{definition}

\begin{theorem}[Erste Orthogonalitätsrelation]
\label{satz:orth1}
Seien $V_1,\ldots,V_r$ die paarweise nicht-isomorphen irreduziblen Darstellungen
einer endlichen Gruppe $G$ über $\mathbb{C}$.  Dann gilt:
\[
    \langle \chi_{V_i}, \chi_{V_j} \rangle_G = \delta_{ij}.
\]
Mit anderen Worten: Die irreduziblen Charaktere bilden eine \emph{Orthonormalbasis}
des Raums $\mathrm{Kl}(G)$ der Klassenfunktionen.
\end{theorem}

\begin{proof}
Aus Schurs Lemma (Korollar~\ref{kor:schur}) folgt:
\[
    \langle \chi_{V_i}, \chi_{V_j} \rangle_G
    = \frac{1}{|G|} \sum_{g \in G} \chi_{V_i}(g)\,\overline{\chi_{V_j}(g)}
    = \dim_{\mathbb{C}} \mathrm{Hom}_G(V_i, V_j)
    = \delta_{ij}. \qedhere
\]
\end{proof}

\begin{theorem}[Zweite Orthogonalitätsrelation]
\label{satz:orth2}
Seien $C_1,\ldots,C_r$ die Konjugationsklassen von $G$ und $g_s \in C_s$
beliebige Repräsentanten.  Dann gilt:
\[
    \sum_{i=1}^{r} \chi_{V_i}(g_s)\,\overline{\chi_{V_i}(g_t)} = \frac{|G|}{|C_s|}\,\delta_{st}.
\]
\end{theorem}

\begin{korollar}[Anzahl der Irreduziblen]
Die Anzahl der irreduziblen Darstellungen einer endlichen Gruppe ist gleich der
Anzahl ihrer Konjugationsklassen.
\end{korollar}

\begin{korollar}[Zerlegungsformel]
\label{kor:zerlegung}
Für jede Darstellung $(V,\rho)$ mit Charakter $\chi_V$ gilt: Die Multiplizität
der irreduziblen Darstellung $V_i$ in $V$ ist
\[
    m_i = \langle \chi_V, \chi_{V_i} \rangle_G.
\]
\end{korollar}

\begin{beispiel}[Charaktertafel von $S_3$]
Die symmetrische Gruppe $S_3$ hat Ordnung $6$ und drei Konjugationsklassen:
$C_1 = \{e\}$, $C_2 = \{(12),(13),(23)\}$, $C_3 = \{(123),(132)\}$.
Die drei irreduziblen Darstellungen sind: die triviale Darstellung $\mathbf{1}$,
die Vorzeichen-Darstellung $\mathrm{sgn}$ und die Standard-Darstellung
$\mathrm{std}$ der Dimension $2$.

\begin{center}
\begin{tabular}{lccc}
\toprule
 & $e$ & $(12)$ & $(123)$ \\
\midrule
$\chi_{\mathbf{1}}$ & $1$ & $1$ & $1$ \\
$\chi_{\mathrm{sgn}}$ & $1$ & $-1$ & $1$ \\
$\chi_{\mathrm{std}}$ & $2$ & $0$ & $-1$ \\
\bottomrule
\end{tabular}
\end{center}

Man verifiziert $\sum_i (\chi_{V_i}(e))^2 = 1 + 1 + 4 = 6 = |S_3|$, eine
notwendige Identität für eine vollständige Liste der Irreduziblen.
\end{beispiel}

\begin{theorem}[Burnsides Dimensionsformel]
\label{satz:burnside_dim}
Sei $G$ eine endliche Gruppe mit irreduziblen komplexen Darstellungen
$V_1,\ldots,V_r$ der Dimensionen $d_1,\ldots,d_r$.  Dann gilt:
\[
    \sum_{i=1}^{r} d_i^2 = |G|.
\]
\end{theorem}

%% ============================================================
\section{Darstellungen der symmetrischen Gruppen $S_n$}
\label{sec:sn}
%% ============================================================

Die symmetrische Gruppe $S_n$ ist eine der bedeutendsten Gruppen in Mathematik
und Physik, und ihre Darstellungstheorie wird vollständig durch die Kombinatorik
von \emph{Partitionen} und \emph{Young-Tableaux} beschrieben.

\begin{definition}
Eine \emph{Partition} von $n$ ist eine schwach absteigende Folge positiver
ganzer Zahlen $\lambda = (\lambda_1 \geq \lambda_2 \geq \cdots \geq \lambda_k > 0)$
mit $\lambda_1 + \cdots + \lambda_k = n$.  Wir schreiben $\lambda \vdash n$.
\end{definition}

\begin{definition}
Das \emph{Young-Diagramm} von $\lambda \vdash n$ besteht aus $n$ Kästchen in
Zeilen von oben nach unten, mit $\lambda_i$ Kästchen in Zeile $i$.  Ein
\emph{Young-Tableau} der Form $\lambda$ ist eine Füllung der Kästchen mit den
Zahlen $1,\ldots,n$; es heißt \emph{standardisiert}, wenn die Einträge entlang
jeder Zeile von links nach rechts und entlang jeder Spalte von oben nach unten
zunehmen.
\end{definition}

\begin{beispiel}
Die Partition $\lambda = (3,2) \vdash 5$ hat Young-Diagramm mit drei Kästchen in
der ersten und zwei in der zweiten Zeile.  Es gibt genau $5$ standardisierte
Young-Tableaux dieser Form.  Zum Beispiel:
\[
\begin{array}{|c|c|c|}
\hline
1 & 2 & 3 \\
\hline
4 & 5 & \\
\cline{1-2}
\end{array}
\quad\text{und}\quad
\begin{array}{|c|c|c|}
\hline
1 & 2 & 4 \\
\hline
3 & 5 & \\
\cline{1-2}
\end{array}
\quad\text{usw.}
\]
\end{beispiel}

\begin{definition}[Specht-Modul]
Für $\lambda \vdash n$ ist der \emph{Specht-Modul} $S^\lambda$ ein spezifischer
irreduzibler $\mathbb{C}[S_n]$-Modul, der wie folgt konstruiert wird.  Man fixiere
ein Young-Tableau $T$ der Form $\lambda$ und setze
\[
    R(T) = \{\sigma \in S_n : \sigma \text{ erhält jede Zeile von } T\},
\quad
    C(T) = \{\sigma \in S_n : \sigma \text{ erhält jede Spalte von } T\}.
\]
Der \emph{Young-Symmetrisierer} ist
\[
    c_T = \underbrace{\left(\sum_{\sigma \in R(T)} \sigma\right)}_{a_T} \cdot
          \underbrace{\left(\sum_{\tau \in C(T)} \mathrm{sgn}(\tau)\,\tau\right)}_{b_T}.
\]
Dann ist $S^\lambda = \mathbb{C}[S_n] \cdot c_T$ ein irreduzibler
$\mathbb{C}[S_n]$-Modul, bis auf Isomorphismus unabhängig von der Wahl von $T$.
\end{definition}

\begin{theorem}[Klassifikation der Irreduziblen von $S_n$]
\label{satz:sn_klass}
Die Moduln $\{S^\lambda : \lambda \vdash n\}$ bilden eine vollständige Familie
paarweise nicht-isomorpher irreduzibler komplexer Darstellungen von $S_n$.  Die
Dimension von $S^\lambda$ ist gleich der Anzahl standardisierter Young-Tableaux
der Form $\lambda$, gegeben durch die \emph{Hakenlängenformel}:
\[
    \dim S^\lambda = \frac{n!}{\prod_{\square \in \lambda} h(\square)},
\]
wobei $h(\square)$ die \emph{Hakenlänge} am Kästchen $\square$ bezeichnet
(Anzahl der Kästchen direkt rechts davon plus direkt darunter, plus $1$ für
$\square$ selbst).
\end{theorem}

\begin{beispiel}[Frobenius-Formel]
Der Charakter $\chi^\lambda$ des Specht-Moduls $S^\lambda$ kann durch die
Frobenius-Charakterformel berechnet werden.  Für eine Permutation $\sigma \in S_n$
vom Zyklentyp $\mu = (1^{m_1} 2^{m_2} \cdots)$ gilt:
\[
    \chi^\lambda(\sigma) = \langle s_\lambda, p_\mu \rangle,
\]
wobei $s_\lambda$ das Schur-Polynom ist, $p_\mu = \prod_i p_{\mu_i}$ die
Potenzreihen-Symmetriefunktion und das innere Produkt das Hall-Skalarprodukt
auf symmetrischen Funktionen.
\end{beispiel}

\begin{beispiel}[Irreduzible von $S_4$]
Die Partitionen von $4$ sind $(4)$, $(3,1)$, $(2,2)$, $(2,1,1)$, $(1,1,1,1)$
und liefern $5$ irreduzible Darstellungen der Dimensionen $1, 3, 2, 3, 1$.
Man prüft: $1^2 + 3^2 + 2^2 + 3^2 + 1^2 = 24 = 4!$ wie erwartet.
\end{beispiel}

%% ============================================================
\section{Darstellungen von $\mathrm{GL}_n$ und $\mathrm{SL}_n$}
\label{sec:gln}
%% ============================================================

Die Darstellungstheorie kompakter Lie-Gruppen, oder allgemeiner reduktiver
algebraischer Gruppen, ist weit reicher als der endliche Fall.  Wir skizzieren
die Höchstgewicht-Klassifikation für $\mathrm{GL}_n(\mathbb{C})$ und
$\mathrm{SL}_n(\mathbb{C})$.

\begin{definition}
Eine Darstellung $(\rho, V)$ von $\mathrm{GL}_n(\mathbb{C})$ heißt
\emph{polynomiell} (bzw.\ \emph{rational}), wenn jeder Matrixkoeffizient
$g \mapsto \langle v^*, \rho(g)v\rangle$ eine polynomielle (bzw.\ rationale)
Funktion der Matrixeinträge von $g$ ist.
\end{definition}

Die irreduziblen polynomiellen Darstellungen von $\mathrm{GL}_n$ werden durch
\emph{dominante Gewichte} indiziert: Folgen $\lambda = (\lambda_1 \geq \lambda_2
\geq \cdots \geq \lambda_n \geq 0)$ nicht-negativer ganzer Zahlen.  Diese
entsprechen genau den Partitionen $\lambda$ mit höchstens $n$ Teilen.

\begin{theorem}[Höchstgewicht-Klassifikation]
\label{satz:hoechstgewicht}
Zu jedem dominanten Gewicht $\lambda = (\lambda_1 \geq \cdots \geq \lambda_n \geq 0)$
existiert eine bis auf Isomorphismus eindeutige irreduzible polynomielle
Darstellung $V_\lambda$ von $\mathrm{GL}_n(\mathbb{C})$.  Jede irreduzible
polynomielle Darstellung ist von dieser Form.  Die Darstellung $V_\lambda$ ist
genau der Schur-Modul zur Partition $\lambda$.
\end{theorem}

\begin{theorem}[Weylsche Charakterformel]
\label{satz:weyl_charformel}
Sei $G = \mathrm{SL}_n(\mathbb{C})$ mit maximalem Torus $T$ der
Diagonalmatrizen der Determinante $1$.  Für ein dominantes ganzzahliges Gewicht
$\lambda$ ist der Charakter der irreduziblen Darstellung $V_\lambda$, ausgewertet
auf $\mathrm{diag}(t_1,\ldots,t_n) \in T$:
\[
    \chi_{V_\lambda}(t) = \frac{\sum_{w \in W} \epsilon(w)\, t^{w(\lambda+\rho)-\rho}}{\prod_{\alpha > 0}(1 - t^{-\alpha})},
\]
wobei $W = S_n$ die Weyl-Gruppe ist, $\epsilon(w) = \det(w)$, $\rho =
\tfrac{1}{2}\sum_{\alpha > 0}\alpha$ der Weyl-Vektor (halbe Summe der positiven
Wurzeln) und das Produkt über positive Wurzeln $\alpha$ läuft.
\end{theorem}

\begin{beispiel}[Darstellungen von $\mathrm{SL}_2$]
Für $G = \mathrm{SL}_2(\mathbb{C})$ sind die dominanten Gewichte die
nicht-negativen ganzen Zahlen $n \geq 0$.  Die irreduzible Darstellung $V_n$
des Höchstgewichts $n$ hat Dimension $n+1$ und kann als Raum der homogenen
Polynome vom Grad $n$ in zwei Variablen realisiert werden, wobei $\mathrm{SL}_2$
durch Substitution wirkt.  Der Charakter ist
\[
    \chi_{V_n}\!\left(\begin{smallmatrix}t & 0 \\ 0 & t^{-1}\end{smallmatrix}\right)
    = t^n + t^{n-2} + \cdots + t^{-n} = \frac{t^{n+1} - t^{-(n+1)}}{t - t^{-1}}.
\]
Dies sind verkleidete Tschebyscheff-Polynome und stimmen mit den Spin-$n/2$-Dar\-stel\-lun\-gen
von $\mathrm{SU}(2)$ in der Quantenmechanik überein.
\end{beispiel}

\begin{proposition}[Schur-Funktoren]
Für eine Partition $\lambda \vdash d$ mit höchstens $n$ Teilen ordnet der
\emph{Schur-Funktor} $\mathbb{S}^\lambda$ jedem Vektorraum $E$ der Dimension $n$
die Darstellung
\[
    \mathbb{S}^\lambda E \coloneqq (E^{\otimes d} \otimes S^\lambda)^{S_d}
\]
zu, wobei $S^\lambda$ der Specht-Modul von $S_d$ ist.  Dies liefert den
irreduziblen $\mathrm{GL}(E)$-Modul $V_\lambda$ und stellt eine kanonische
Verbindung zwischen den Darstellungstheorien von $S_d$ und $\mathrm{GL}_n$ her.
\end{proposition}

%% ============================================================
\section{Anwendungen}
\label{sec:anwendungen}
%% ============================================================

\subsection{Burnsides $pq$-Satz}

\begin{theorem}[Burnside, 1904]
\label{satz:burnside_pq}
Sei $G$ eine endliche Gruppe der Ordnung $p^a q^b$, wobei $p$ und $q$ Primzahlen
und $a, b \geq 0$ sind.  Dann ist $G$ auflösbar.
\end{theorem}

\begin{proof}[Beweisidee]
Der Beweis verwendet zwei charaktertheoretische Hilfssätze.

\begin{lemma}
\label{lem:b1}
Ist $\chi$ ein irreduzibler Charakter von $G$ vom Grad $d$ und $C$ eine
Konjugationsklasse mit $\gcd(|C|, d) = 1$, so gilt für jedes $g \in C$ entweder
$\chi(g) = 0$ oder $\rho(g)$ ist eine skalare Matrix.
\end{lemma}

\begin{lemma}
\label{lem:b2}
Eine endliche Gruppe ist genau dann auflösbar, wenn jeder nicht-triviale
irreduzible Charakter einen Grad hat, der durch eine Primzahl teilbar ist.
\end{lemma}

Mittels Sylow-Theorie findet man in $G$ eine Konjugationsklasse $C$ der Größe
$p^s$ für ein $s \geq 1$.  Lemma~\ref{lem:b1} erzwingt, dass jeder irreduzible
Charakter vom Grad nicht durch $p$ teilbar entweder auf $C$ verschwindet oder
$\rho(g)$ skalar macht.  Die Spalten-Orthogonalitätsrelation erzwingt dann eine
normale Untergruppe, was die Einfachheit von $G$ widerlegt, außer wenn $G$ abelsch
ist.  Induktion über $|G|$ liefert die Auflösbarkeit.
\end{proof}

\begin{bemerkung}
Burnsides ursprünglicher Beweis von 1904 war der erste Einsatz von Charaktertheorie
zum Beweis eines rein gruppentheoretischen Resultats.  Ein charakterfreier Beweis
wurde von Goldschmidt und Bender erst in den 1970er Jahren gefunden und ist
erheblich aufwändiger.  Dies illustriert die Stärke darstellungstheoretischer
Methoden.
\end{bemerkung}

\subsection{Molekulare Symmetrie}

In der Chemie wirkt die Symmetriegruppe eines Moleküls (wie die Rotationsgruppe
eines regulären Polyeders) auf dem Raum der Elektronenorbitale oder
Schwingungsmoden.  Die Zerlegung dieser (typischerweise reduzibler) Darstellungen
in Irreduzible erklärt die Entartungsstruktur der Energieniveaus und die
Auswahlregeln für spektroskopische Übergänge.

Beispielsweise hat das Methan-Molekül $\mathrm{CH}_4$ die Symmetriegruppe
$T_d \cong S_4$.  Die $2p$-Orbitale der vier Wasserstoffatome bilden die
$4$-dimensionale Permutationsdarstellung, die als $A_1 \oplus T_2$ zerfällt
(in spektroskopischer Notation) und das beobachtete Bindungsmuster erklärt.

\subsection{Kristallographie}

Die $230$ kristallographischen Raumgruppen in drei Dimensionen klassifizieren alle
möglichen periodischen Atomanordnungen.  Jede solche Gruppe $\Gamma$ passt in eine
kurze exakte Sequenz $1 \to \mathbb{Z}^3 \to \Gamma \to P \to 1$, wobei $P$ eine
\emph{Punktgruppe} (eine endliche Untergruppe von $\mathrm{O}(3)$) ist.  Die
irreduziblen Darstellungen von $P$ klassifizieren die Symmetrietypen der
Phonon-Moden am $\Gamma$-Punkt der Brillouin-Zone, mit direkten Anwendungen
in der Interpretation von Raman- und Infrarotspektren.

Darüber hinaus wird das \emph{Burnside-Lemma} (Cauchy-Frobenius-Lemma)
\[
    |X/G| = \frac{1}{|G|} \sum_{g \in G} |X^g|
\]
in der Kombinatorik eingesetzt, um die Anzahl wesentlich verschiedener Objekte
unter Symmetriewirkung zu zählen --- z.B.\ die Anzahl chemisch verschiedener
Färbungen eines Moleküls.

%% ============================================================
\section{Verbindung zum Langlands-Programm}
\label{sec:langlands}
%% ============================================================

Das Langlands-Programm, das Robert Langlands 1967 in einem berühmten Brief an
André Weil initiierte, ist ein weitreichendes Geflecht von Vermutungen und Sätzen,
das tiefe Korrespondenzen zwischen drei scheinbar unabhängigen Objekten vorhersagt:
\begin{enumerate}
  \item \textbf{Automorphe Darstellungen}: bestimmte Darstellungen adelischer
        Gruppen $G(\mathbb{A}_{\mathbb{Q}})$, die Modulformen und Dirichlet-Charaktere
        verallgemeinern.
  \item \textbf{Galois-Darstellungen}: stetige Homomorphismen $\rho \colon
        \mathrm{Gal}(\bar{\mathbb{Q}}/\mathbb{Q}) \to {}^L G$ in die
        Langlands-duale Gruppe.
  \item \textbf{L-Funktionen}: meromorphe Funktionen auf $\mathbb{C}$, die
        arithmetische Information in Euler-Produkten codieren und
        Funktionalgleichungen erfüllen.
\end{enumerate}

\subsection{Galois-Darstellungen}

Sei $\ell$ eine Primzahl.  Eine \emph{$\ell$-adische Galois-Darstellung} ist
ein stetiger Homomorphismus
\[
    \rho \colon \mathrm{Gal}(\bar{\mathbb{Q}}/\mathbb{Q}) \longrightarrow \mathrm{GL}_n(\mathbb{Z}_\ell),
\]
wobei $\mathbb{Z}_\ell$ die $\ell$-adische Topologie und die Galois-Gruppe die
Krull-Topologie trägt.  Die in den vorigen Abschnitten entwickelte
Darstellungstheorie gilt hier mit $G = \mathrm{Gal}(\bar{\mathbb{Q}}/\mathbb{Q})$,
einer profinierten Gruppe.

\begin{beispiel}[Zyklotomischer Charakter]
Der \emph{zyklotomische Charakter} $\chi_\ell \colon
\mathrm{Gal}(\bar{\mathbb{Q}}/\mathbb{Q}) \to \mathbb{Z}_\ell^\times$ ist durch
$\sigma(\zeta_{\ell^n}) = \zeta_{\ell^n}^{\chi_\ell(\sigma)}$ für alle
$\ell^n$-ten Einheitswurzeln $\zeta_{\ell^n}$ definiert.  Dies ist eine
irreduzible eindimensionale Galois-Darstellung, die die Arithmetik
zyklotomischer Körper kodiert.
\end{beispiel}

\begin{beispiel}[Galois-Darstellung einer elliptischen Kurve]
Sei $E/\mathbb{Q}$ eine elliptische Kurve.  Das Tate-Modul
\[
    T_\ell(E) = \varprojlim_n E[\ell^n](\bar{\mathbb{Q}}) \cong \mathbb{Z}_\ell^2
\]
trägt eine natürliche Wirkung von $\mathrm{Gal}(\bar{\mathbb{Q}}/\mathbb{Q})$
und liefert eine zweidimensionale $\ell$-adische Darstellung
$\rho_{E,\ell} \colon \mathrm{Gal}(\bar{\mathbb{Q}}/\mathbb{Q}) \to \mathrm{GL}_2(\mathbb{Z}_\ell)$.
Die L-Funktion $L(E,s) = L(\rho_{E,\ell}, s)$ ist dieselbe L-Funktion, die in
der Birch-und-Swinnerton-Dyer-Vermutung auftritt.
\end{beispiel}

\subsection{Die Langlands-Korrespondenz}

\begin{vermutung}[Globale Langlands-Korrespondenz für $\mathrm{GL}_n$]
\label{verm:langlands}
Es gibt eine natürliche Bijektion zwischen:
\begin{enumerate}[label=\textup{(\roman*)}]
  \item Kuspidalen automorphen Darstellungen $\pi$ von $\mathrm{GL}_n(\mathbb{A}_{\mathbb{Q}})$.
  \item Irreduziblen $n$-dimensionalen komplexen Darstellungen der
        Weil--Deligne-Gruppe $W_{\mathbb{Q}}' = W_{\mathbb{Q}} \ltimes \mathbb{C}$
        (oder, vermutungsweise, der Langlands-Gruppe $\mathcal{L}_{\mathbb{Q}}$).
\end{enumerate}
Unter dieser Bijektion stimmt die L-Funktion $L(s,\pi)$ auf der automorphen Seite
mit der Artin-L-Funktion $L(s,\rho)$ auf der Galois-Seite überein.
\end{vermutung}

\begin{bemerkung}
Vermutung~\ref{verm:langlands} ist im Allgemeinen offen.  Wesentliche
Spezialfälle wurden bewiesen:
\begin{itemize}
  \item $n = 1$: Klassenkörpertheorie (Kronecker--Weber, Artin-Reziprozität).
  \item $G = \mathrm{GL}_2$, Gewicht-$2$-Formen: Taylor--Wiles (1995), als Teil
        des Beweises des Letzten Satzes von Fermat.
  \item Funktionskörper-Analogon: Drinfeld (1974, $n=2$) und Lafforgue
        (2002, allgemeines $n$).
  \item Geometrisches Langlands-Programm über komplexen Kurven: Frenkel--Ben-Zvi
        u.a.
\end{itemize}
Die jüngsten Arbeiten von V.\ Lafforgue über die Richtung ``automorph zu Galois''
über Funktionskörpern sowie von Scholze--Fargues mittels perfektoider Geometrie
und der Fargues--Fontaine-Kurve stellen die aktuelle Forschungsfront dar.
\end{bemerkung}

\subsection{Modularität und Fermats Letzter Satz}

Die Verbindung zwischen Darstellungstheorie und Zahlentheorie wird besonders
sichtbar im Beweis des Letzten Satzes von Fermat.  Angenommen, es gäbe eine
Lösung $a^p + b^p = c^p$ mit $p \geq 5$ prim.  Frey (1986) ordnete dieser
hypothetischen Lösung die elliptische Kurve $E_{a,b,c}$ zu, und Ribet (1990)
bewies, dass die zugehörige Galois-Darstellung $\rho_{E,p}$ eine
Gewicht-$2$-Modulform des Niveaus $2$ liefern würde.  Da es jedoch keine solchen
Formen gibt (durch Aufzählung kleiner Niveaus), ergibt sich ein Widerspruch.
Der entscheidende Schritt---dass jede semistabile elliptische Kurve über
$\mathbb{Q}$ modular ist (d.h.\ ihre Galois-Darstellung von einer Modulform
stammt)---wurde von Wiles und Taylor--Wiles 1995 bewiesen.  Dieser
Modularitätssatz ist ein Meilenstein der Langlands-Korrespondenz in Dimension
$n = 2$.

%% ============================================================
\section*{Literatur}
%% ============================================================

\begin{thebibliography}{99}

\bibitem{Serre77}
J.-P.\ Serre,
\textit{Lineare Darstellungen endlicher Gruppen}
[Originalausgabe: \textit{Linear Representations of Finite Groups}],
Graduate Texts in Mathematics, Bd.~42,
Springer-Verlag, New York, 1977.

\bibitem{FultonHarris91}
W.\ Fulton und J.\ Harris,
\textit{Representation Theory: A First Course},
Graduate Texts in Mathematics, Bd.~129,
Springer-Verlag, New York, 1991.

\bibitem{JamesLiebeck01}
G.\ James und M.\ Liebeck,
\textit{Representations and Characters of Groups},
2.\ Aufl., Cambridge University Press, Cambridge, 2001.

\bibitem{BumpLie04}
D.\ Bump,
\textit{Lie Groups},
Graduate Texts in Mathematics, Bd.~225,
Springer-Verlag, New York, 2004.

\bibitem{Langlands70}
R.\ P.\ Langlands,
\textit{Problems in the Theory of Automorphic Forms},
in: Lectures in Modern Analysis and Applications III,
Lecture Notes in Mathematics, Bd.~170,
Springer-Verlag, Berlin, 1970, S.~18--61.

\bibitem{TaylorWiles95}
R.\ Taylor und A.\ Wiles,
Ring theoretic properties of certain Hecke algebras,
\textit{Ann.\ of Math.} \textbf{141} (1995), 553--572.

\bibitem{Lafforgue02}
L.\ Lafforgue,
Chtoucas de Drinfeld et correspondance de Langlands,
\textit{Invent.\ Math.} \textbf{147} (2002), 1--241.

\bibitem{Neukirch99}
J.\ Neukirch, A.\ Schmidt und K.\ Wingberg,
\textit{Cohomology of Number Fields},
Grundlehren der mathematischen Wissenschaften, Bd.~323,
Springer-Verlag, Berlin, 1999.

\bibitem{BumpSchilling17}
D.\ Bump und A.\ Schilling,
\textit{Crystal Bases: Representations and Combinatorics},
World Scientific, Singapur, 2017.

\bibitem{Bourbaki68}
N.\ Bourbaki,
\textit{Groupes et algèbres de Lie}, Chapitres 4--6,
Hermann, Paris, 1968.

\end{thebibliography}

\end{document}
